\documentclass{article}
\usepackage{graphicx}
\usepackage{booktabs}
\usepackage[top=2cm, bottom=2cm, left=2.5cm, right=2.5cm]{geometry}
\usepackage{microtype}

\begin{document}

\begin{center}
  {\large\bfseries COMS4033A/7044A: RBC Report}\\[4pt]
  Dylan Davies \quad 2672598 \\[2pt]
  Sudhir Ajodhya  \quad 2652330
\end{center}

% Intro
\section{Introduction}

Reconnaissance Blind Chess (RBC) is a variant of chess in which each player observes only
their own pieces; the opponent's pieces are hidden until sensed. Agents must therefore
maintain a \emph{belief state} — a set of board positions consistent with all observations
so far — and make decisions under uncertainty. This report presents the development and
evaluation of two agents: \textit{RandomSensing} (baseline) and \textit{ImprovedAgent}.
Multiple improvements were explored and implemented to create a competitive agent capable
of intelligent sensing and cautious tactical play.

% Results
\section{Tournament Results}

The double round-robin tournament involved four agents, with each agent playing twice
against every opponent (once as White, once as Black) for a total of 12 games.
Table~\ref{tab:leaderboard} details the results.

\begin{table}[h!]
\centering
\caption{Agent matchup leaderboard}
\label{tab:leaderboard}
\begin{tabular}{lccccccc}
\toprule
\textbf{Implementation} & \textbf{W} & \textbf{L} & \textbf{D} & \textbf{E} & \textbf{Win\%} & \textbf{As White\%} & \textbf{As Black\%} \\
\midrule
ImprovedAgent & 6 & 0 & 0 & 0 & 100.0\% & 100.0\% & 100.0\% \\
RandomSensing & 4 & 2 & 0 & 0 & 66.7\%  & 66.7\%  & 66.7\%  \\
TroutBot      & 2 & 4 & 0 & 0 & 33.3\%  & 33.3\%  & 33.3\%  \\
RandomBot     & 0 & 6 & 0 & 0 & 0.0\%   & 0.0\%   & 0.0\%   \\
\bottomrule
\end{tabular}
\end{table}

ImprovedAgent achieved a perfect 6--0 record, winning every game as both White and
Black. Its informed sensing strategy kept the belief state small and accurate, enabling
Stockfish to operate on reliable boards and produce consistently strong moves.
RandomSensing finished second at 4--2, defeating RandomBot and TroutBot but losing
both games to ImprovedAgent, confirming that even uninformed sensing with majority
voting outperforms purely random play. TroutBot placed third at 2--4; its Stockfish
engine crashed in several games, forcing it to forfeit those encounters and explaining
its underperformance relative to its expected strength. RandomBot lost all six games,
as expected for a purely random mover.

% Improvements
\section{Improvements to ImprovedAgent}

\subsection{Information-Theoretic Sensing}

The most significant improvement replaced random sensing with a disagreement-based
strategy. For each candidate sense centre, the agent computes a window score by
summing a per-square contribution over all 9 squares of the $3{\times}3$ window:

\[
  \text{score}(sq) = |S| - \max_{p} \bigl|\{s \in S : \text{piece}(s,\, sq) = p\}\bigr|
\]

This measures how many boards would be eliminated if the most-common piece hypothesis
at that square were wrong — i.e.\ how much the square contributes to disambiguation.
The window with the highest total score is selected; ties are broken by proximity to
the board centre. Own pieces are excluded since their positions are always known. If
all scores are zero the agent falls back to a random interior square.

For the first two turns the agent senses the board centre (e4 as White, d5 as Black)
to gather broad early information before the belief set has been meaningfully narrowed.

\subsection{Weighted Move Scoring}

Rather than a simple vote count, each Stockfish suggestion is assigned a weighted score
before aggregation:

\begin{itemize}
  \item \textbf{Base weight:} 1.0 per recommendation.
  \item \textbf{+0.7} if the move is a capture on the current board.
  \item \textbf{+1.0} if the move gives check on the current board.
\end{itemize}

Captures and checks therefore accumulate higher totals and are preferred when raw vote
counts are tied, nudging the agent toward tactically decisive play.

\subsection{King-Capture Confidence Threshold and Belief State Cap}

RandomSensing attempts a king capture whenever it receives the most votes, even if
those votes represent a small fraction of all tracked boards, risking wasted turns
chasing phantom kings. ImprovedAgent only commits to a king capture if the top
candidate exceeds $20\%$ of all possible boards, balancing reliability and
aggression.

The belief state cap for opponent-move expansion was raised to $20{,}000$ states.
Because the improved sensing strategy eliminates far more inconsistent boards per
turn than random sensing does, the true belief state converges faster; a larger cap
therefore adds relatively little computation while allowing a more faithful
representation of the true board distribution.

\subsection{Move-Result Filtering}

After each move, the agent applies four consistency filters to prune the belief state:
(i)~boards on which a requested-but-blocked move was actually legal are discarded;
(ii)~boards on which the move taken was illegal are discarded;
(iii)~if a capture occurred, boards that would not have produced that capture are
discarded; (iv)~if no capture occurred, boards that would have captured are discarded.
This tightens the belief state after every action, compounding the gains from
informed sensing.

% Conclusion
\section{Conclusion}

ImprovedAgent's perfect tournament record demonstrates that principled uncertainty
management is the dominant factor in RBC performance. The information-theoretic
sensing strategy was the primary driver: by directing the sense window toward squares
of maximum disagreement, the belief state is pruned aggressively each turn, giving
Stockfish accurate boards to evaluate. Weighted move scoring and the king-capture
confidence threshold further reduce costly errors, while post-move filtering compounds
the belief-state gains after every action. Together, these improvements produce an
agent that outperforms all baselines regardless of colour.

\end{document}
